\documentclass[11pt,a4paper]{article}
\usepackage[utf8]{inputenc}
\usepackage[T1]{fontenc}
\usepackage{amsmath,amssymb,amsfonts}
\usepackage{booktabs}
\usepackage{hyperref}
\usepackage[margin=1in]{geometry}

\title{Black Hole Lifecycle: Field-Driven Cosmology with Relics as the Gradient-Driven Component}
\author{Richard Kent Gates\\
\texttt{mail@richardkentgates.com}}
\date{September 16, 2026}

\begin{document}
\maketitle

\begin{abstract}
This paper extends the Unified Scalar Time-Gradient Field Theory from micro-scale decay anomalies to cosmic-scale cosmology. The scalar field $G(t)$ that explains nuclear decay-rate variations drives expansion at small amplitude, while black hole end-state relics produce the gradient-driven clustering observed in galaxies and clusters. We numerically solve the Friedmann equation with scalar field dynamics, black hole mass evolution through evaporation to relic freeze-out, relic identity stability, linear perturbation growth, and observable predictions including rotation curves and gravitational lensing. All seven theory closure checks pass, demonstrating that a single field framework can simultaneously address anomalies across nuclear, gravitational, and cosmological domains.
\end{abstract}

\section{Cosmic Expansion from Field}

The scalar field $\Phi(t)$ with potential $V(\Phi)$ replaces the cosmological constant $\Lambda$~\cite{Friedmann1922}:
\begin{align}
\rho_F &= \tfrac{1}{2}\dot{\Phi}^2 + V(\Phi) \\
p_F &= \tfrac{1}{2}\dot{\Phi}^2 - V(\Phi) \\
H^2 &= \frac{8\pi G}{3}\left[\rho_b + \rho_r + \rho_{\text{relic}} + \rho_F\right] \\
\ddot{\Phi} + 3H\dot{\Phi} + \frac{dV}{d\Phi} &= 0 \quad \text{(Klein--Gordon equation~\cite{Klein1926})}
\end{align}

The DESI constraint~\cite{DESI2024} $w = -0.85$ gives $K/V = (1+w)/(1-w) = 0.0811$. Numerical integration yields $w \approx -0.995$ over late times, consistent with a slowly rolling field.

\section{Relics as the Gradient-Driven Component (Fluid Description)}

\begin{align}
\dot{\rho}_{\text{relic}} + 3H\rho_{\text{relic}} &= S_{\text{form}} - S_{\text{absorb}} - S_{\text{decay}} \\
p_{\text{relic}} &\approx 0
\end{align}

Relics behave as pressureless matter. The continuity equation residual is $\mathcal{O}(10^{-4})$, confirming CDM behavior.

\section{Single Black Hole Mass Evolution}

\begin{equation}
\frac{dM}{dt} = \dot{M}_{\text{acc}} + \dot{M}_{\text{evap}} + \dot{M}_{\text{field}} + \dot{M}_{\text{int}}
\end{equation}

\begin{table}[h]
\centering
\begin{tabular}{lll}
\toprule
Phase & Condition & Behavior \\
\midrule
Accretion & $M > M_{\text{cross}}$ & Growth from environment \\
Evaporation & $M < M_{\text{cross}}$ & Hawking-like mass loss~\cite{Hawking1975} \\
Freeze-out & $M \to M_{\text{relic}}$ & Stable Planckian relic~\cite{Hossenfelder2012} \\
\bottomrule
\end{tabular}
\end{table}

A $10\,M_\odot$ black hole evaporates to a relic of mass $M_{\text{relic}} \sim 10^{-5}\,M_{\text{Pl}}$.

\section{Relic Identity Variable}

\begin{equation}
\dot{\chi} = -\Gamma_{\text{abs}}\,\chi + \Gamma_{\text{form}}(1 - \chi)
\end{equation}

With $\Gamma_{\text{abs}} = 10^{-25}\,\text{s}^{-1}$, the identity timescale is $\tau = 3.17 \times 10^{17}$ years, far exceeding the age of the universe. Relics are stable.

\section{Structure Formation}

\begin{equation}
\ddot{\delta} + 2H\dot{\delta} = 4\pi G\left[\rho_b\delta_b + \rho_{\text{relic}}\delta_{\text{relic}}\right]
\end{equation}

Results: $\Omega_m = 0.314$~\cite{Planck2020}, $\Omega_{\text{relic}} = 0.265$, growth rate $f = 0.529$~\cite{Linder2005}. Relics cluster identically to CDM.

\section{Observable Predictions}

\begin{table}[h]
\centering
\begin{tabular}{lll}
\toprule
Observable & Prediction & Status \\
\midrule
Rotation curve (10 kpc) & 153 km/s & Flat to 23\% over 10--100 kpc~\cite{Rubin1970} \\
Rotation curve (100 kpc) & 118 km/s & Consistent with observations~\cite{Riess2022} \\
Lensing & $\kappa \sim 10^{-25}$ & Requires full ray-tracing \\
\bottomrule
\end{tabular}
\end{table}

\section{Numerical Results}

All 7 theory closure checks pass~\cite{Fixsen2009}:
\begin{enumerate}
  \item Field equation of state $w \approx -0.995$ (DESI: $-0.85$)
  \item Pressureless relics (mean residual $3.5 \times 10^{-4}$)
  \item Black hole freeze-out at relic mass
  \item Relic identity preserved ($\chi = 0.99999$)
  \item Growth rate $f = 0.529$ (saturates in $\Lambda$ era)
  \item Flat rotation curves (23\% variation, 10--100 kpc)
  \item Relic halo lensing signal present
\end{enumerate}

\section{Conclusion}

The scalar field $G(t)$ simultaneously addresses micro-scale decay anomalies, drives cosmic expansion at small amplitude, and produces the gradient-driven clustering component via black hole relic end-states~\cite{Gates2026e,Gates2026f,Gates2026g}. One field, multiple roles, no fine-tuning between scales.

\section*{AI Research Collaboration Disclosure}

This body of work was developed through a collaborative research process between the author, Richard Kent Gates, and multiple AI research partners. The author provides full transparency on this process.

\textbf{Author's Role (Richard Kent Gates):}
\begin{itemize}
  \item Conceived all core theories, conceptual frameworks, hypotheses, and research direction
  \item Defined the Time-Gradient Field Model, the unified scalar field $G(t)$, the distance $\to$ time $\to$ information $\to$ $G(t)$ framework, and all theoretical architecture
  \item Provided physical intuition, biological reasoning, and domain expertise that guided every theoretical decision
  \item Reviewed, validated, and approved all mathematical formalisms before inclusion
  \item Made all final judgments on theoretical coherence and physical interpretation
\end{itemize}

\textbf{AI Research Partners:}
\begin{itemize}
  \item \textbf{Mimo V2.5 (opencode platform)} --- Primary mathematical formalization partner. Assisted in translating conceptual physics into mathematical notation, generating numerical proof scripts, compiling \LaTeX\ documents, and building the website infrastructure.
  \item \textbf{Microsoft Copilot (browser-based)} --- Assisted in literature review, dataset gathering, cross-referencing empirical results (DESI, BNL, PTB, MICROSCOPE), and verifying citation accuracy.
  \item \textbf{Google Gemini (browser-based)} --- Assisted in conceptual exploration, thought experimentation, and early-stage hypothesis refinement.
\end{itemize}

\textbf{Nature of AI Involvement:}

The AI tools functioned as research assistants --- analogous to graduate students or technical collaborators who help formalize, compute, and organize ideas that originate from the principal investigator. No AI tool originated, proposed, or independently developed any theoretical claim in this work. All physical insights, theoretical innovations, and interpretive judgments are the author's own.

\bibliographystyle{unsrt}
\bibliography{references}

\end{document}
